% Mathématiques 4e — cours V3
% Proportionnalité et quatrième proportionnelle

\section{Objectifs}

\begin{itemize}
\item Reconnaître une situation de proportionnalité.
\item Utiliser un coefficient de proportionnalité.
\item Déterminer une quatrième proportionnelle.
\item Choisir entre linéarité, retour à l'unité, coefficient et produit en croix.
\item Calculer et appliquer un pourcentage.
\item Utiliser des échelles, agrandissements et réductions simples.
\item Représenter une relation de proportionnalité par un tableau ou un graphique.
\item Reconnaître graphiquement une proportionnalité.
\item Résoudre des problèmes contextualisés faisant intervenir plusieurs grandeurs.
\end{itemize}

\section{1. Deux grandeurs proportionnelles}

Deux grandeurs sont proportionnelles si l'on peut passer de l'une à l'autre en multipliant toujours par le même nombre.

Ce nombre est le coefficient de proportionnalité.

Si :
\[
y=kx,
\]
alors $k$ est ce coefficient.

\begin{exemple}
Un article coûte :
\[
2{,}40\ \text{€}
\]
l'unité.

Pour $x$ articles, le prix vaut :
\[
P=2{,}40x.
\]

Le coefficient est :
\[
2{,}40.
\]
\end{exemple}

\section{2. Reconnaître une situation}

Une situation n'est pas proportionnelle simplement parce que deux grandeurs varient ensemble.

Si une location coûte :
\[
5\ \text{€}
\]
de frais fixes puis :
\[
3\ \text{€}
\]
par heure, le prix n'est pas proportionnel à la durée.

Pour une heure :
\[
8\ \text{€}.
\]

Pour deux heures :
\[
11\ \text{€}.
\]

Doubler la durée ne double pas le prix.

\section{3. Tableau de proportionnalité}

Un tableau peut organiser les couples de valeurs.

\begin{tabular}{c|c|c|c}
Quantité & $3$ & $5$ & $8$ \\
Prix & $12$ & $20$ & $32$ \\
\end{tabular}

Le coefficient permettant de passer de la première ligne à la seconde vaut :
\[
4.
\]

\coursefigure{/images/mathematiques/4eme/proportionnalite-tableau.svg}{Tableau de proportionnalité illustrant la recherche d'une quatrième valeur.}{Le même coefficient relie toutes les colonnes d'un tableau de proportionnalité.}

\section{4. Linéarité multiplicative}

Si une valeur est multipliée par un nombre, sa correspondante est multipliée par le même nombre.

Par exemple :
\[
3\longrightarrow12.
\]

Comme :
\[
9=3\times3,
\]
alors :
\[
9\longrightarrow12\times3=36.
\]

Cette méthode est particulièrement rapide lorsque le facteur est simple.

\section{5. Linéarité additive}

On peut additionner deux correspondances.

Si :
\[
3\longrightarrow12
\]
et :
\[
5\longrightarrow20,
\]
alors :
\[
8=3+5
\]
correspond à :
\[
12+20=32.
\]

\section{6. Retour à l'unité}

Si :
\[
7\ \text{kg}
\]
coûtent :
\[
19{,}60\ \text{€},
\]
le prix d'un kilogramme vaut :
\[
19{,}60\div7=2{,}80\ \text{€}.
\]

Pour :
\[
11\ \text{kg},
\]
on obtient :
\[
11\times2{,}80=30{,}80\ \text{€}.
\]

\section{7. Quatrième proportionnelle}

On connaît trois nombres dans un tableau de proportionnalité et on cherche le quatrième.

\[
\begin{array}{c|cc}
& 3 & 8\\
\hline
&12&x
\end{array}
\]

Ici, le coefficient est :
\[
12\div3=4.
\]

Donc :
\[
x=8\times4=32.
\]

Le nombre :
\[
32
\]
est une quatrième proportionnelle.

\section{8. Produit en croix}

En 4e, on peut utiliser l'égalité des produits en croix dans une situation de proportionnalité.

Si :
\[
\dfrac ab=\dfrac cd,
\]
alors :
\[
ad=bc.
\]

Pour :
\[
\dfrac38=\dfrac{x}{40},
\]
on écrit :
\[
8x=3\times40.
\]

Donc :
\[
8x=120
\]
et :
\[
x=15.
\]

\begin{attention}
Le produit en croix est une conséquence d'une égalité de quotients. Il ne doit pas être appliqué mécaniquement à un tableau qui n'est pas proportionnel.
\end{attention}

\section{9. Choisir la méthode}

Une quatrième proportionnelle peut être obtenue de plusieurs façons.

\begin{methode}
\begin{itemize}
\item facteur évident : linéarité multiplicative ;
\item combinaison simple de colonnes : linéarité additive ;
\item valeur unitaire utile : retour à l'unité ;
\item nombres moins commodes : coefficient ou produit en croix.
\end{itemize}
\end{methode}

La méthode la plus courte n'est pas toujours la plus mécanique.

\section{10. Pourcentage}

Un pourcentage est une proportion sur :
\[
100.
\]

\[
p\%=\dfrac{p}{100}.
\]

Pour calculer :
\[
18\%
\]
de :
\[
250,
\]
on écrit :
\[
\dfrac{18}{100}\times250=45.
\]

Donc :
\[
\boxed{18\%\text{ de }250=45}.
\]

\section{11. Retrouver un pourcentage}

Dans un groupe de :
\[
80
\]
personnes, :
\[
28
\]
répondent oui.

La proportion vaut :
\[
\dfrac{28}{80}=0{,}35.
\]

Donc :
\[
\boxed{35\%}
\]
ont répondu oui.

\section{12. Augmentation en pourcentage}

Un prix de :
\[
120\ \text{€}
\]
augmente de :
\[
15\%.
\]

Le montant de l'augmentation vaut :
\[
0{,}15\times120=18\ \text{€}.
\]

Le nouveau prix est :
\[
120+18=138\ \text{€}.
\]

\begin{remarque}
Une augmentation de $15\%$ ne signifie pas ajouter le nombre $15$.
\end{remarque}

\section{13. Réduction en pourcentage}

Un article de :
\[
80\ \text{€}
\]
bénéficie d'une réduction de :
\[
25\%.
\]

La réduction vaut :
\[
0{,}25\times80=20\ \text{€}.
\]

Le prix final vaut :
\[
80-20=60\ \text{€}.
\]

\section{14. Échelle}

À l'échelle :
\[
1:50\,000,
\]
un centimètre sur la carte représente :
\[
50\,000\ \text{cm}
\]
dans la réalité.

Or :
\[
50\,000\ \text{cm}=500\ \text{m}.
\]

Donc :
\[
1\ \text{cm}\longrightarrow500\ \text{m}.
\]

\section{15. Exemple d'échelle}

Deux points sont séparés de :
\[
7{,}4\ \text{cm}
\]
sur une carte au :
\[
1:25\,000.
\]

La distance réelle vaut :
\[
7{,}4\times25\,000=185\,000\ \text{cm}.
\]

Donc :
\[
185\,000\ \text{cm}=1{,}85\ \text{km}.
\]

\section{16. Agrandissement et réduction}

Une maquette réalisée à l'échelle :
\[
1:20
\]
est une réduction.

Une longueur réelle de :
\[
3\ \text{m}=300\ \text{cm}
\]
devient :
\[
300\div20=15\ \text{cm}.
\]

Le coefficient de réduction vaut :
\[
\dfrac1{20}.
\]

\section{17. Graphique}

Une relation proportionnelle donne des points alignés avec l'origine du repère.

\coursefigure{/images/mathematiques/4eme/proportionnalite-graphique.svg}{Comparaison entre une relation proportionnelle et une droite ne passant pas par l'origine.}{L'alignement seul ne suffit pas : la droite associée à une proportionnalité doit passer par l'origine.}

\begin{propriete}
Une relation de proportionnalité est représentée graphiquement par une droite passant par :
\[
O(0;0).
\]
\end{propriete}

\section{18. Exemple graphique}

Les points :
\[
(1;3),\quad(2;6),\quad(4;12),\quad(5;15)
\]
sont alignés avec l'origine.

Le rapport :
\[
\dfrac yx=3
\]
est constant.

La relation est :
\[
y=3x.
\]

\section{19. Relation affine non proportionnelle}

Les points de :
\[
y=3x+4
\]
sont alignés.

Mais pour :
\[
x=0,
\]
on obtient :
\[
y=4.
\]

La droite ne passe pas par l'origine.

La relation n'est donc pas proportionnelle.

\section{20. Vitesse constante}

Un véhicule roule à :
\[
72\ \text{km/h}
\]
à vitesse constante.

En :
\[
2{,}5\ \text{h},
\]
la distance vaut :
\[
72\times2{,}5=180\ \text{km}.
\]

La distance est proportionnelle à la durée si la vitesse est constante.

\section{21. Débit constant}

Un robinet remplit :
\[
18\ \text{L}
\]
en :
\[
3\ \text{min}.
\]

Le débit vaut :
\[
18\div3=6\ \text{L/min}.
\]

En :
\[
7\ \text{min},
\]
le volume vaut :
\[
6\times7=42\ \text{L}.
\]

\section{22. Proportionnalité en géométrie}

Un agrandissement de coefficient :
\[
1{,}5
\]
multiplie toutes les longueurs par :
\[
1{,}5.
\]

Une longueur de :
\[
8\ \text{cm}
\]
devient :
\[
12\ \text{cm}.
\]

Le même coefficient doit s'appliquer à toutes les longueurs correspondantes.

\section{23. Exemple développé}

Une photographie mesure :
\[
12\ \text{cm}\times18\ \text{cm}.
\]

On l'agrandit en conservant les proportions pour obtenir une largeur de :
\[
30\ \text{cm}.
\]

Le coefficient est :
\[
30\div18=\dfrac53.
\]

La nouvelle hauteur vaut :
\[
12\times\dfrac53=20\ \text{cm}.
\]

Le format agrandi est :
\[
\boxed{20\ \text{cm}\times30\ \text{cm}}.
\]

\section{24. Méthode générale}

\begin{methode}
\begin{enumerate}
\item identifier les deux grandeurs et leurs unités ;
\item vérifier que la situation est proportionnelle ;
\item organiser les valeurs dans un tableau si utile ;
\item choisir la méthode adaptée ;
\item calculer la valeur cherchée ;
\item contrôler le coefficient et l'ordre de grandeur ;
\item interpréter avec l'unité.
\end{enumerate}
\end{methode}

\section{25. Erreurs fréquentes}

\begin{itemize}
\item Utiliser un produit en croix sans vérifier la proportionnalité.
\item Confondre coefficient de proportionnalité et valeur recherchée.
\item Ajouter un pourcentage comme un nombre absolu.
\item Mélanger des unités dans une échelle.
\item Croire que des points simplement alignés suffisent graphiquement.
\item Oublier le passage par l'origine.
\item Appliquer un coefficient différent à deux longueurs correspondantes.
\end{itemize}

\section{À retenir}

\begin{aretenir}
Dans une situation proportionnelle :
\[
\boxed{y=kx}.
\]

Pour chercher une quatrième proportionnelle, on peut utiliser la linéarité, le retour à l'unité, le coefficient ou l'égalité des produits en croix.

Graphiquement :
\[
\boxed{\text{les points sont alignés avec l'origine}}.
\]

Un pourcentage est une proportion sur $100$ et une échelle est une relation de proportionnalité entre longueurs correspondantes.
\end{aretenir}
